\section{Outpacing DRAM Refresh} \label{sec:fast}
    To achieve "\textit{\hyperref[ch:fast]{C3. Outpacing DRAM Refresh}}", memory accesses must race against the next 
    row refresh in order to induce 
    sufficiently frequent activation's of the same \gls{dram} rows. Accordingly, this section focuses on 
    the implementation of access patterns and parameter configurations for high-frequency activation 
    with some optimization techniques, and adopts concrete techniques for handling and bypassing 
    mitigation mechanisms to sustain effective hammering on ARM-based systems such as Raspberry Pi 
    devices.

    \subsection{TRR Detection Approaches} \label{subsec:trr-detection-approaches}
        \gls{trr} is a standard hardware protection mechanism designed to successfully 
        prevent both single-sided and double-sided Rowhammer attacks~\cite{paperTRRScope, paperTRRespass}. When a Rowhammer attack 
        targets a specific row, \gls{trr} responds by generating an additional refresh operation to refresh 
        the corresponding victim row, thereby mitigating disturbance effects~\cite{webTenYearsOfRH}. The primary objective 
        of this analysis is to determine whether \gls{trr} is available and active on the test systems used in 
        this work.
        \\
        To investigate this, the official JEDEC documentation for LPDDR4~\cite{docJEDECLPDDR4} and LPDDR4x~\cite{docJEDECLPDDR4x} was 
        examined first. This analysis revealed that the section describing \gls{trr} mode had been removed 
        from the current official documentation. Nevertheless, it was noted that the Mode Register 0 
        configuration specifies that \gls{trr} support is indicated when \texttt{MR0 OP[2]} is set to \texttt{0b}. In 
        parallel, the Raspberry Pi devices were physically inspected by unscrewing them in order to 
        identify the \gls{dram} manufacturer, the component markings, and the FBGA numbers. Using the “FBGA 
        to component marking decoder” provided on the official manufacturer website~\cite{webFBGADecoder}, the exact 
        part numbers of the \gls{dram} components were determined, as summarized in~\cref{tab:fbga-part-number-decoder}.
        
        \begin{table}[h!]
            \centering
            \caption{FBGA Part Number Decoder}
            \label{tab:fbga-part-number-decoder}
            \begin{tabular}{lllll}
                \hline
                \textbf{Model} & \textbf{RAM} & \textbf{Manufacturer} & \textbf{FBGA} & \textbf{Part Number}\\
                \hline
                Raspberry Pi 4 & 4 GB & Micron & D9WHV & MT53D1024M32D4DT-053 WT:D \\
                Raspberry Pi 4 & 8 GB & Micron & D9ZCL & MT53E2G32D4NQ-046 WT:A \\
                Raspberry Pi 5 & 4 GB & Micron & D8CJG & MT53E1G32D2FW-046 WT:B \\
                Raspberry Pi 5 & 8 GB & Micron & D8CJN & MT53E2G32D4DE-046 WT:C \\
                Raspberry Pi 5 & 16 GB & Micron & D8CBG & MT53E4G32D8CY-046 WT:C \\
                \hline
            \end{tabular}
        \end{table}
        
        The identification of the exact \gls{dram} components enabled conclusions to be drawn regarding the 
        officially applicable documentation for the test devices~\cite{datasheetHammerpi02, datasheetHammerpi03, datasheetHammerpi04, datasheetHammerpi05}. However, this 
        documentation could not be accessed, even after contacting the manufacturers directly, as no 
        response was received. Despite this limitation, an older Micron reference document from 2022 
        was identified~\cite{datasheetLPDDR4}, which provides sufficient information about the specific devices used 
        in this work. Importantly, this reference still includes the \gls{trr} Mode section that has been 
        removed from newer documentation.
        \\
        From this reference, the theoretical foundations of \gls{trr} were derived, including the concept 
        of the \gls{mac} and the system's response once this limit is reached. The 
        \gls{mac} specifies the maximum number of activations that a single row, denoted as a target row 
        \texttt{TRn}, can sustain within a refresh period before adjacent rows, defined as victim rows, 
        must be refreshed to prevent disturbance errors~\cite{datasheetLPDDR4}. When the \gls{mac} limit is reached, the device 
        either performs all required refresh commands before allowing further row activations or enters 
        \gls{trr} mode, in which only the adjacent victim rows are refreshed~\cite{datasheetLPDDR4}.
        \\
        An analysis of the relevant mode register configuration, as summarized in~\cref{tab:mode-register-assignments}, 
        reveals that \texttt{MR0} is used for mode selection, controlling either RFM or 
        or \gls{trr} operation. Specifically, \gls{trr} control applies only when \texttt{MR0 OP[2] = 0b}, whereas RFM is 
        selected when \texttt{MR0 OP[2] = 1b}. The \gls{trr} mode itself is defined by \texttt{MR24 OP[7]}, where 
        a value of \texttt{1} enables \gls{trr} mode and \texttt{0} disables it. The \gls{trr} mode BAn, which is 
        indicating the bank of the target row, is specified by \texttt{MR24 OP[6:4]}. Furthermore, 
        \texttt{MR24 OP[3]} determines the \gls{mac} type, with \texttt{0} corresponding to a limited \gls{mac} and 
        \texttt{1} to an unlimited \gls{mac}. The \gls{mac} value is encoded in \texttt{MR24 OP[2:0]}, where \texttt{000b} 
        represents an unlimited \gls{mac}, while any other configuration specifies a finite \gls{mac} value. Notably, 
        when an unlimited \gls{mac} is configured by setting \texttt{MR24 OP[2:0] = 000b} and 
        \texttt{MR24 OP[3] = 1}, \gls{trr} operation is not required, even if \gls{trr} mode itself is enabled via
        \texttt{MR24 OP[7]}. By default, \gls{trr} mode is disabled, as indicated by \texttt{MR24 OP[7] = 0b}. 
        Despite this, the exact device behavior is vendor-specific.

        \begin{table}[h!]
            \centering
            \caption{Mode Register Assignments}
            \label{tab:mode-register-assignments}
            \resizebox{\textwidth}{!}{%
                \begin{tabular}{|l|l|l|l|l|l|l|l|l|l|}
                    \hline
                    \textbf{MR\#} & \textbf{Function} & \textbf{OP7} & \textbf{OP6} & \textbf{OP5} & \textbf{OP4} & \textbf{OP3} & \textbf{OP2} & \textbf{OP1} & \textbf{OP0} \\
                    \hline
                    0 & Mode Select & \multicolumn{2}{|l|}{} & Singleended mode & \multicolumn{2}{l|}{} & XORed RFM/TRR support & Latency mode & REF \\
                    \hline
                    24 & TRR control & TRR mode & \multicolumn{3}{|l|}{TRR mode BAn} & UnItd MAC & \multicolumn{3}{|l|}{MAC value} \\
                    \hline
                    24 & RFM control & \multicolumn{2}{|l|}{RAAMMT} & \multicolumn{5}{|l|}{RAAIMT} & RFM \\
                    \hline
                \end{tabular}%
            }
        \end{table}
        
        To confirm or refute whether \gls{trr} is implemented on the test devices, a deeper inspection of the 
        relevant mode register entries is required, specifically \texttt{MR0 OP[2]}, \texttt{MR24 OP[3:0]}, 
        and \texttt{MR24 OP[7]}.
        \\
        Since mode registers can only be accessed through the memory controller, an 
        indirect approach was considered. The \gls{spd} mechanism was identified as a 
        potential source of information, as it stores memory-related parameters such as \gls{mac} and possibly \gls{trr} 
        configuration~\cite{webSPD}. While DIMM-based systems typically include an \gls{spd} EEPROM, it is unclear whether 
        such an EEPROM exists for soldered \gls{dram}, as used in Raspberry Pi devices. The proposed approach was 
        therefore to access the system bus using I2C tools, dump the \gls{spd} EEPROM if present, and decode its 
        contents according to the corresponding \gls{ddr} standard.
        \\
        The workflow began with the initial detection of available I2C buses using \texttt{sudo i2cdetect -l}, 
        which returned no entries, and \texttt{sudo i2cdetect -y 1}, which failed with an error indicating that 
        \texttt{/dev/i2c-1} was not found. Consequently, the I2C interface was enabled via \texttt{sudo raspi-config}
        by selecting within the Interface Options the enabling I2C option. 
        After reboot, it was assumed that the I2C driver would load during boot and that the \texttt{/dev/i2c-*} 
        device files would become available.
        \\
        Subsequent verification using \texttt{ls /dev/i2c*} confirmed the presence of three I2C buses: 
        \texttt{/dev/i2c-1}, \texttt{/dev/i2c-13}, and \texttt{/dev/i2c-14}. Listing all I2C adapters with 
        \texttt{sudo i2cdetect -l} showed one standard Synopsys DesignWare adapter and two SoC-integrated 
        adapters. Unfortunately, a scan of the standard I2C bus using 
        \texttt{sudo i2cdetect -y 1} did not detect any devices across the standard address range used by 
        \gls{spd} EEPROMs (e.g., 0x50-0x57), as all addresses were reported as unused.
        \\
        These results confirm that the I2C bus is active and operational, but no external I2C devices, 
        including an \gls{spd} EEPROM at standard addresses, were detected. This observation is consistent with 
        the fact that Raspberry Pi devices use soldered \gls{dram} and aligns with previous findings reported in~\cite{paperProFlipping}.
        Consequently, \texttt{decode-dimms}, which is a primary x86 tool for reading memory device 
        information from DIMM EEPROMs, could not be used. Furthermore, no interface exists for directly 
        reading device information from the \gls{dram} mode registers. Although it could theoretically be assumed 
        that \gls{trr} is not enabled by default according to~\cite{datasheetLPDDR4}, this behavior is vendor-specific. Therefore, 
        without access to non-public vendor documentation, the default \gls{trr} configuration of the test devices 
        cannot be conclusively determined.

    \subsection{Bypass Techniques} \label{subsec:bypass-techniques}
        Given the inability to conclusively determine whether \gls{trr} is implemented on the test devices, \gls{trr} 
        must be assumed to be present. This assumption is further supported by the fact that previous 
        Rowhammer tests failed to induce bit flips, even though the official Micron documentation suggests 
        that \gls{trr} is disabled by default~\cite{datasheetLPDDR4}. To address this, established \gls{trr} bypass techniques were 
        considered, including Rowhammer fuzzers such as Blacksmith~\cite{toolBlacksmith} and TRRespass~\cite{toolTRRespass}.
        \\
        Blacksmith~\cite{paperBlacksmith} is a scalable Rowhammer fuzzer primarily designed for x86 systems. Introduced 
        in 2022, it constructs novel non-uniform hammering patterns based on frequency, phase, and 
        amplitude to bypass modern in-\gls{dram} \gls{trr} mitigations and successfully induce bit flips. TRRespass~\cite{paperTRRespass},
        introduced in 2020, is another Rowhammer fuzzer for x86 systems that automatically 
        identifies \gls{trr}-bypassing access patterns. It generates a wide range of effective Rowhammer patterns 
        on \gls{trr}-protected \gls{dram} modules by varying two parameters, namely cardinality and location, and by 
        employing non-conventional, many-sided Rowhammer patterns to circumvent both memory-controller-level 
        and \gls{dram}-level defenses. Within this work, the focus is placed on porting TRRespass, as the 
        authors of~\cite{paperTRRespass} implemented a simplified ARMv8 version to test LPDDR4 and LPDDR4x memory and 
        study the prevalence of the issue, although no public repository for this implementation exists.
        \\
        Both Blacksmith and TRRespass rely on the availability of HugePages~\cite{webHugePagesRPi4, toolMemoryAware, webHugepagesDebian}. Initially, 
        HugePages could not be used because the running kernel did not support HugeTLB or Transparent 
        HugePages, as indicated by the absence of any output from \texttt{grep Huge /proc/meminfo}. An 
        inspection of the kernel configuration using \texttt{/boot/config-\$(uname -r)} revealed that 
        while the architecture fundamentally supports HugePages, the kernel was compiled without 
        enabling either Transparent HugePages or HugeTLBFS. As a result, HugePages were enabled by 
        natively rebuilding the Linux kernel, as described in~\cref{sec:lab-setup}.
        \\
        According to the Debian arm64 kernel documentation, HugeTLB page sizes of 2 MB and 1 GB are 
        supported, but 1 GB HugeTLB pages must be pre-allocated at boot time via kernel command line 
        arguments, as they cannot be allocated dynamically~\cite{webHugepagesDebian}. Only smaller page sizes, such as 
        64 KB, 2 MB, and 32 MB, can be allocated at runtime. Therefore, the kernel command line was 
        extended to pre-allocate a single 1 GB HugeTLB page.
        \\
        After enabling HugePages, we verified their availability. On the 
        Raspberry Pi 4, the available sizes were 64 KB, 2 MB, 32 MB, and 1 GB, whereas on the 
        Raspberry Pi 5 only 2 MB, 32 MB, and 1 GB HugePages were supported. \\
        To ensure deterministic 
        HugeTLB behavior, Transparent HugePages were explicitly disabled via a privileged write to 
        \texttt{/sys/kernel/mm/transparent\_hugepage/enabled}. Following kernel command line 
        configuration and a system reboot, the kernel message buffer was filtered using 
        \texttt{dmesg | grep -i huge} to verify which HugePage sizes were successfully pre-allocated. 
        Finally, the active HugePage configuration was confirmed by querying 
        \texttt{/sys/kernel/mm/hugepages/*/nr\_hugepages}, which reports the number of HugePages 
        currently in use for each supported size.
    
    \subsection{TRRespass ARM Port} \label{subsec:trrespass-arm-port}
        This subsection focuses exclusively on the architectural adaptations required to port TRRespass from 
        x86-based systems to the ARMv8 architecture in order to enable execution on the ARM-based 
        single-board computers used in this work, denoted as test devices.
        \\
        The scope of this porting effort is deliberately constrained. The core fuzzing and hammering 
        algorithms implemented in TRRespass remain unchanged, and no modifications are made to their 
        fundamental logic or operation. Consequently, the algorithmic behavior and the effectiveness 
        of the underlying attacks are not reassessed in this work, as these aspects are thoroughly 
        analyzed and discussed in the original TRRespass publication~\cite{paperTRRespass}.
        \\
        Several porting challenges arise from the fundamental differences between x86 and ARMv8 
        architectures. A primary issue is the absence of x86-specific instructions and registers on 
        ARM platforms, which necessitates the replacement of x86-only CPU instructions with 
        ARMv8-compatible equivalents. This challenge is further compounded by the lack of direct 
        counterparts for certain cache flush and timing instructions that are commonly available on 
        x86 systems. In addition, compatibility with Linux hugepage support on ARM systems must be
        ensured, as outlined in the previous section. The migration also requires replacing x86-specific 
        compiler flags with configurations suitable for ARMv8. To support heterogeneous ARM platforms 
        across different test devices, system page sizes must be handled dynamically. Furthermore, 
        hugepages need to be integrated to preserve contiguous memory regions, which are essential for 
        both hammering and fuzzing operations.
        \\
        As a result of these adaptations, the modified TRRespass tool executes successfully on the 
        ARMv8-based platforms evaluated in~\cite{paperBachelorarbeit}. Moreover, the ported TRRespass 
        tool~\cite{toolHamOnRaspi} establishes a foundation for future research on Rowhammer-related attacks and 
        defenses in ARM-based systems.
